We developed a hybrid hidden Markov model that discretizes equity excess growth rates into Laplace quantile-defined states and augments regime-switching dynamics with a Poisson-driven jump-duration mechanism. Applied to SPY over a ten-year training window and validated out-of-sample on the full calendar year 2025, the framework achieves KS and AD pass rates exceeding 95\% while reproducing the volatility clustering that standard regime-switching models miss. The jump mechanism, governed by only two scalar hyperparameters, is the structural feature responsible for generating persistent high-volatility regimes, and the direct frequentist estimation strategy eliminates the computational cost and initialisation sensitivity of the EM algorithm. A Single-Index Model extension scales the approach to a 424-asset universe, providing a baseline for multi-asset synthetic scenario generation.

Several extensions follow naturally. Time-varying transition matrices, estimated via rolling windows or Bayesian online learning, would allow the model to adapt to structural regime shifts. Replacing the single-factor SIM with multi-factor specifications such as the Fama-French three- or five-factor model \cite{fama_french_1993}, principal-component-based factor structures, or nonlinear residual generators could substantially improve asset-level distributional fidelity. Adaptive state definitions based on density clustering rather than fixed quantiles could improve tail coverage during crisis periods. Finally, embedding the hybrid HMM as the generative component within a portfolio optimisation loop, where synthetic scenarios drive tail-risk objectives, would provide an end-to-end application of the framework.
